\SourceLeaf{284}%
\begin{center}\begin{tabular}{@{}r c@{}}
8.1 & $-$\\
9.2 & $+$\\
12.5 & $-$\\
14.7 & $+$\\
16.11 & $-$\\
17.4 & $-$\\
18.3 & $+$\\
19.10 & $+$\\
20.15 & $+$\\
21.6 & $-$\\
22.13 & $-$\\
\end{tabular}\end{center}

Clearly, if we start from a point before (1) and traverse
points 1, 2, 3 etc. along the curve, a figure first forms at
crossing 8. Let us call it $a$. It runs from 1
to 8. Evidently the part of the whole line immediately before 1 and
[the part] after 8 lie on the $-$ side of $a$.

Continuing in this way, a new connected
curve forms that runs from the starting point to 1 (then leaving the previous path 2, 3 etc.)
and goes on to 9, 10, 11 etc. At 9 one reaches the
$+$ side of $a$, at 12 the $-$ side, and at 14 the $+$ side of $a$.
At 16, however, the curve meets itself and closes a second
figure, which runs from 11 through 12, 13, 14, 15 to 16. Call it $b$.
Evidently the part of the curve immediately before 1, then the entire
piece from 8 to 11, and finally the part immediately adjoining 16 all lie on
the $-$ side of $b$. Now that $b$ has also been separated, the
curve meets $a$ again at 17 and comes to the $-$ side, and meets it again at 18,
coming to the $+$ side. At 19 the $3^{\mathrm{rd}}$
figure closes (since 10 does not yet lie in an excluded part). Call it
$c$. It runs from 10 to 11 $=$ 16, then from 16 to 19. And
when this figure too has been separated, the piece left over
and the piece that immediately follows lie on the $+$ side of $c$. At 20 one reaches
the $+$ side of $b$, at 21 the $-$ side of $a$, and at 22
the $-$ side of $b$. As one returns to the starting point (before 1),
the $4^{\mathrm{th}}$ figure, which we call $d$,
closes.
\par\Needspace{30\baselineskip}
The position relative to the curves already traced is thus as follows
